\section{Preliminaries: the anatomy of headset-to-arm teleoperation}
\label{sec:prelim}

This section introduces the parts every headset-driven teleoperation
system shares, so that a reader new to the area can follow the design
choices of Sections~\ref{sec:system} and~\ref{sec:method}. We cover the
common data flow (Section~\ref{sec:prelim-dataflow}), the retargeting
maps that turn hand poses into end-effector targets
(Section~\ref{sec:prelim-retarget}), the inverse-kinematics families
that turn targets into joint angles (Section~\ref{sec:prelim-ik}), the
standard input filter (Section~\ref{sec:prelim-filter}), and the common
treatments of workspace limits (Section~\ref{sec:prelim-workspace}).
Throughout, the discussion applies to single or dual arms with five,
six, or seven degrees of freedom; we note where the arm's DoF count
changes the picture.

\subsection{Data flow}
\label{sec:prelim-dataflow}

A headset such as the Meta Quest tracks its two hand controllers at
\SIrange{50}{90}{Hz} and reports, per controller, a 6-D pose (position
and orientation), analogue grip and trigger values, and button states.
The poses arrive in a \emph{tracking frame} whose origin and yaw depend
on where the headset started~--- an arbitrary frame the system must not
control in directly. A teleoperation pipeline therefore runs the
following stages, usually as separate threads at different rates:
controller poses are (i) re-expressed in a stable \emph{operator frame},
(ii) low-pass filtered, (iii) mapped by a \emph{retargeting map} to
end-effector pose targets in each arm's \emph{base frame}, (iv) solved
by inverse kinematics into joint targets, and (v) streamed to the servo
controllers, typically at \SI{100}{Hz} or more. The analogue grips
usually gate the whole pipeline (the \emph{clutch},
Section~\ref{sec:prelim-retarget}) and the triggers drive the grippers.
We write \(p_H, R_H\) for the filtered hand position and orientation in
the operator frame, and \(p_E, R_E\) for the end-effector pose in the
arm's base frame.

\subsection{Retargeting maps}
\label{sec:prelim-retarget}

\paragraph{Position: absolute vs clutched.}
The simplest map is absolute: fix a rigid transform between the operator
frame and the robot frame once, and command
\(p^\ast = T\,p_H\) forever after. Absolute mapping preserves a
persistent hand--robot correspondence but fails when the robot workspace
and the comfortable human workspace differ in size or offset~--- which
is the usual case. The standard remedy is the \emph{clutched} (or
indexed) map: an engage/disengage control, usually the grip buttons,
gates tracking. On engagement the system latches the current hand
position \(p_H^0\) and end-effector position \(p_E^0\), and while
engaged it commands the anchored delta
\begin{equation}
  p^\ast(t) \;=\; p_E^0 + s\,\bigl(p_H(t) - p_H^0\bigr),
  \label{eq:prelim-clutch}
\end{equation}
with a translation scale \(s\). Disengaging freezes the robot;
re-engaging re-anchors. The operator can therefore cover any workspace
in comfortable strokes, like lifting and re-placing a computer mouse.

\paragraph{Orientation.}
For arms with a full 3-DoF wrist (6- and 7-DoF arms), orientation is
usually retargeted absolutely: the target rotation is the hand rotation,
optionally composed with a fixed offset, and a 7th degree of freedom is
spent on posture objectives such as elbow placement. At five degrees of
freedom the picture changes: with only two wrist joints (flex and roll),
the end-effector yaw is kinematically coupled to the shoulder-pan angle,
so the reachable orientation set at a given position is a 2-D subset of
\(\mathrm{SO}(3)\). Commanding an arbitrary hand rotation then
over-constrains the tracker. Systems escape in one of two ways: they
\emph{relax} the tracker (lower or zero the orientation cost on the
infeasible axis and accept whatever orientation results), or they
\emph{construct} targets inside the reachable set before tracking. The
armplane mapping of Section~\ref{sec:armplane} is of the second kind.

\subsection{Inverse-kinematics families}
\label{sec:prelim-ik}

Four families cover the systems in this paper.
\emph{Differential (velocity-level) IK} linearises the kinematics each
tick and solves a quadratic programme for joint velocities that reduce
task errors under joint position and velocity limits; frame-tracking,
damping, and posture objectives enter as weighted tasks. It is smooth
and constraint-aware, and its per-tick cost is low.
\emph{Damped least squares} is the classical closed-form member of the
same family: it applies the damped pseudo-inverse of the task Jacobian
to the pose error, trading tracking gain against robustness near
singularities~\citep{wampler1986,nakamura1986}.
\emph{Per-tick nonlinear optimisation} solves the position-level problem
directly each tick (bounded least squares over joint angles). It can be
the most accurate, but each solve is expensive and consecutive solves
may jump between IK branches, so a rate limiter downstream is essential.
\emph{Analytic and split solutions} exploit the arm's specific geometry:
low-DoF chains often admit closed-form wrist solutions, so a system can
solve position with a small numeric problem and set the wrist joints
directly. Split methods buy wrist agility at the price of coupling
between the two subproblems.

\subsection{Input filtering}
\label{sec:prelim-filter}

Tracked controller poses carry hand tremor and sensor noise. Filtering
them with a fixed low-pass forces a hard choice between residual jitter
and lag. The One-Euro filter~\citep{oneeuro}, the de-facto standard for
interactive input, is a first-order low-pass whose cutoff frequency
adapts to the signal's speed:
\(f_c = f_{c,\min} + \beta\,|\dot{\hat x}|\), where the speed estimate
\(\dot{\hat x}\) is itself low-pass filtered with cutoff \(f_{c,d}\).
The two knobs separate the concerns: \(f_{c,\min}\) sets steadiness when
the hand is still, and \(\beta\) sets how quickly the filter opens up
during fast motion. Position and orientation are filtered separately;
quaternions need a sign-consistency guard so that interpolation stays on
the short arc.

\subsection{Workspace handling}
\label{sec:prelim-workspace}

Every arm has a bounded reachable set, and headset teleoperation
routinely commands targets outside it. With no handling, the IK layer
chases an infeasible target: the QP saturates against joint limits, the
arm parks at a limit posture away from its commanded target, and
re-entry is abrupt. Common remedies operate on the target before the IK
sees it: \emph{projection} replaces an outside target with the nearest
point of a reachability approximation, \emph{freezing} holds the last
feasible target until the hand returns, and \emph{velocity shaping}
attenuates the outward component of target motion as the boundary
approaches. All three need a cheap membership test, which is why
practical systems use analytic envelope approximations rather than exact
reachability maps. Section~\ref{sec:envelope} makes these concrete for
the SO-101 and compares them head-to-head.
